% Part: normal-modal-logic
% Chapter: completeness
% Section: completeness-K

\documentclass[../../../include/open-logic-section]{subfiles}

\begin{document}

\olfileid[fr]{nml}{com}{cmk}

\olsection{Détermination et complétude pour \Log{K}}

Nous pouvons maintenant utiliser le modèle canonique pour établir la
complétude. Celle-ci découle du fait que les !!{formula}s vraies dans
le modèle canonique de~$\Sigma$ sont exactement celles qui sont
!!{derivable}s dans~$\Sigma$. On dit qu'un modèle ayant cette propriété
\emph{détermine}~$\Sigma$.

\begin{defn}
  Un modèle $\mModel{M}$ \emph{détermine} une logique modale
  normale~$\Sigma$ si et seulement si $\mSat{M}{!A}$ équivaut à $\Sigma
  \Proves !A$, pour toute !!{formula}~$!A$.
\end{defn}

\begin{thm}[Détermination]\ollabel{thm:determination}
  Pour une logique modale normale et cohérente~$\Sigma$, $\mSat{M^\Sigma}{!A}$ si et seulement si $\Sigma \Proves !A$.
\end{thm}

\begin{proof}
  Si $\mSat{M^\Sigma}{!A}$, alors, pour tout ensemble $\Sigma$-cohérent
  complet $\Delta$, on a $\mSat{M^\Sigma}{!A}[\Delta]$. Par le
  lemme de vérité, $!A \in \Delta$ pour tout ensemble $\Sigma$-cohérent
  complet $\Delta$. D'après le
  \olref[lin]{cor:provability-characterization} (avec $\Gamma =
  \emptyset$), on en déduit $\Sigma \Proves !A$.

  Réciproquement, si $\Sigma \Proves !A$, alors, d'après la
  \olref[ccs]{prop:ccs-properties}\olref[ccs]{prop:ccs-closed}, tout
  ensemble $\Sigma$-cohérent complet $\Delta$ contient~$!A$. Par le
  lemme de vérité, $\mSat{M^\Sigma}{!A}[\Delta]$ pour tout~$\Delta \in
  W^\Sigma$ ; autrement dit, $\mSat{M^\Sigma}{!A}$.
\end{proof}

L'absurde $\lfalse$ est faux dans tout modèle. Par correction,
il n'est pas dérivable dans~\Log{K}, qui est donc cohérente.
Le modèle canonique de~\Log{K} détermine alors~\Log{K}. On en déduit
le résultat suivant.

\begin{cor}\ollabel{cor:Kcomplete}
  La logique modale de base \Log{K} est complète relativement à la
  classe de tous les modèles : si $\Entails !A$, alors $\Log{K}
  \Proves !A$.
\end{cor}

\begin{proof}
  Par contraposition, si $\Log{K} \Proves/ !A$, le théorème de
  détermination donne $\mSat/{M^{\Log{K}}}{!A}$ ; donc $!A$ n'est pas
  valide.
\end{proof}

Pour établir par cette méthode la complétude d'un système~$\Sigma$
relativement à une classe de modèles, on vérifie en outre que son
modèle canonique appartient à la classe. Pour \Log{KTB4}, par exemple,
on cherche à établir que la relation canonique est réflexive,
symétrique et transitive. L'appartenance à la classe demande une
preuve et peut échouer.
\end{document}
